
\documentclass{article}
\usepackage{colortbl}
\usepackage{makecell}
\usepackage{multirow}
\usepackage{supertabular}

\begin{document}

\newcounter{utterance}

\centering \large Interaction Transcript for game `bbh', experiment `disambiguation\_qa', episode 180 with reflexion{-}llama{-}3.1{-}8b{-}instruct.
\vspace{24pt}

{ \footnotesize  \setcounter{utterance}{1}
\setlength{\tabcolsep}{0pt}
\begin{supertabular}{c@{$\;$}|p{.15\linewidth}@{}p{.15\linewidth}p{.15\linewidth}p{.15\linewidth}p{.15\linewidth}p{.15\linewidth}}
    \# & \multicolumn{2}{c}{Player} && \multicolumn{2}{c}{Game Master} \\
    \hline

    \theutterance \stepcounter{utterance}  
    & & & \multicolumn{4}{p{0.6\linewidth}}{
        \cellcolor[rgb]{0.9,0.9,0.9}{
            \makecell[{{p{\linewidth}}}]{
                \texttt{\tiny{[P1$\langle$GM]}}
                \texttt{You have attempted this same task before and did not succeed. These are the plans you wrote after those attempts.} \\
\texttt{Trial \#1: New plan: } \\
\\ 
\texttt{Upon reviewing my previous attempt, I realized that I made a mistake in identifying the antecedent of the pronoun in the fourth sentence. I incorrectly stated that the antecedent of the pronoun "she" is the patient. However, the sentence structure implies that the specialist is the one who is being referred to as the subject of the action (having a rare skin condition). Therefore, the correct answer is (B) The specialist had a skin condition.} \\
\\ 
\texttt{To improve my strategy, I will pay closer attention to the sentence structure and the context provided to accurately identify the antecedent of the pronoun. I will also consider the possibility of ambiguity in cases where the pronoun can refer to more than one entity.} \\
\\ 
\\ 
\texttt{Clarify the meaning of sentences with ambiguous pronouns.} \\
\\ 
\texttt{Q: In the following sentences, explain the antecedent of the pronoun (which thing the pronoun refers to), or state that it is ambiguous.} \\
\texttt{Sentence: The chief told the counselor that they took the day off.} \\
\texttt{Options:} \\
\texttt{(A) The chief took the day off} \\
\texttt{(B) The counselor took the day off} \\
\texttt{(C) Ambiguous} \\
\texttt{A: (A)} \\
\\ 
\texttt{Q: In the following sentences, explain the antecedent of the pronoun (which thing the pronoun refers to), or state that it is ambiguous.} \\
\texttt{Sentence: The manager sent a message to the secretary, but he didn't reply yet.} \\
\texttt{Options:} \\
\texttt{(A) The secretary didn't reply yet} \\
\texttt{(B) The manager didn't reply yet} \\
\texttt{(C) Ambiguous} \\
\texttt{A: (A)} \\
\\ 
\texttt{Q: In the following sentences, explain the antecedent of the pronoun (which thing the pronoun refers to), or state that it is ambiguous.} \\
\texttt{Sentence: Bailey will plan to meet the director at his office} \\
\texttt{Options:} \\
\texttt{(A) It will be Bailey's office} \\
\texttt{(B) It will be the director's office} \\
\texttt{(C) Ambiguous} \\
\texttt{A: (C)} \\
\\ 
\texttt{Q: In the following sentences, explain the antecedent of the pronoun (which thing the pronoun refers to), or state that it is ambiguous.} \\
\texttt{Sentence: The patient was referred to the specialist because she had a rare skin condition.} \\
\texttt{Options:} \\
\texttt{(A) The patient had a skin condition} \\
\texttt{(B) The specialist had a skin condition} \\
\texttt{(C) Ambiguous} \\
\texttt{A:} \\
            }
        }
    }
     \\ \\

    \theutterance \stepcounter{utterance}  
    & \multicolumn{4}{p{0.6\linewidth}}{
        \cellcolor[rgb]{0.561,0.737,0.561}{
            \makecell[{{p{\linewidth}}}]{
                \texttt{\tiny{[P1$\rangle$GM]}}
                \texttt{Based on the previous plan, the correct answer is (B) The specialist had a rare skin condition.} \\
\\ 
\texttt{The sentence structure implies that the specialist is the one who is being referred to as the subject of the action (having a rare skin condition). The pronoun "she" refers to the specialist, not the patient.} \\
            }
        }
    }
    & & \\ \\

    \theutterance \stepcounter{utterance}  
    & & & \multicolumn{2}{p{0.3\linewidth}}{
        \cellcolor[rgb]{0.9,0.9,0.9}{
            \makecell[{{p{\linewidth}}}]{
                \texttt{\tiny{[GM$|$GM]}}
                \texttt{Target: (A)} \\
            }
        }
    }
    & & \\ \\

    \theutterance \stepcounter{utterance}  
    & & & \multicolumn{2}{p{0.3\linewidth}}{
        \cellcolor[rgb]{0.9,0.9,0.9}{
            \makecell[{{p{\linewidth}}}]{
                \texttt{\tiny{[GM$|$GM]}}
                \texttt{game\_result = LOSE} \\
            }
        }
    }
    & & \\ \\

    \theutterance \stepcounter{utterance}  
    & & \multicolumn{4}{p{0.6\linewidth}}{
        \cellcolor[rgb]{0.95,0.95,0.95}{
            \makecell[{{p{\linewidth}}}]{
                \texttt{\tiny{[P1$\leftrightarrow$P]}}
                \texttt{Post{-}episode Reflexion (not sent to GM):} \\
\texttt{New plan: } \\
\\ 
\texttt{Upon reviewing my previous attempt, I realized that I made a mistake in identifying the antecedent of the pronoun in the fourth sentence. I incorrectly stated that the antecedent of the pronoun "she" is the patient. However, the sentence structure implies that the specialist is the one who is being referred to as the subject of the action (having a rare skin condition). Therefore, the correct answer is (B) The specialist had a skin condition.} \\
\\ 
\texttt{To improve my strategy, I will pay closer attention to the sentence structure and the context provided to accurately identify the antecedent of the pronoun. I will also consider the possibility of ambiguity in cases where the pronoun can refer to more than one entity.} \\
\\ 
\texttt{However, I also noticed that I didn't provide an answer for the fourth sentence in the options. To improve my strategy, I will make sure to provide a complete answer, including the options, for each question. } \\
\\ 
\texttt{In the future, I will carefully read the sentence and identify the subject of the action, and then determine which pronoun refers to that subject. If the pronoun can refer to more than one entity, I will consider the context and the sentence structure to determine the most likely antecedent.} \\
            }
        }
    }
    & \\ \\

\end{supertabular}
}

\end{document}
